% == == == == == == == == == == == == == == == == == == == == == == == == == ==
    == == == == == == == == == == == ==
    = % Section 3b : Functional Block Diagram % == == == == == == == == == == ==
      == == == == == == == == == == == == == == == == == == == == == == == == ==
      == ==
    =

\section {
  Functional Block Diagram
}
\sectionauthors { All }

\subsection{Functional Overview}

A Functional Block
Diagram(FBD)
provides a detailed view of the system's functional components, showing how different subsystems interact and communicate. While the Conceptual Block Diagram (CBD) presents a high-level overview, the FBD expands on this by illustrating the specific hardware components, interfaces, and data flows within the system.

    The FBD in Figure \ref{fig : fbd} details the STAR Robot's architecture, including the Raspberry Pi 5 as the main compute unit running ROS 2, the ESP32-S3 microcontroller handling low-level motor control and temperature sensing, the various sensors (LiDAR, IMU, ultrasonic), and the power management system. This diagram clarifies how data flows between components and which interfaces (SPI, I\textsuperscript{2}C, GPIO, SMBus) connect the different subsystems.

\begin{figure}[H]
\centering
\fbox {
  \includegraphics[width = 0.95\textwidth] { figures / star_fbd.pdf }
}
\caption { System Functional Block Diagram }
\label {
fig:
  fbd
}
\end { figure }
